%\subsection{A lower bound of $\Omega(n^{1/4}\alpha^{1/2})$ for testing $C_4$-freeness}\label{sec:C4-lb}

\subsection{A two-sided error lower bound for constant-arboricity graphs}
We start by proving the lower bound for constant-arboricity graphs, as stated in Theorem~\ref{thm:C4-5} in the introduction.

\begin{theorem} \label{thm:C4-const-alpha-lb}
Testing $C_4$-freeness in constant-arboricity graphs over $n$ vertices requires $\Omega(n^{1/4})$ queries for constant $\epsilon$.
\end{theorem}
We note that Theorem~\ref{thm:C4-const-alpha-lb} holds for two-sided error algorithms and not only one-sided error ones.
% \commentt{ Give result with dependence on $\eps$?} NO..

\begin{proof}
For any (sufficiently large) $n$ we shall define two bipartite graphs, $G_0$ and $G_1$, where both have constant arboricity. The graph $G_0$ is $C_4$-free and  the graph $G_1$ is $\Omega(1)$-far from being $C_4$ free.
Based on each $G_b$, $b\in \bitset$ we have a distribution, $\calD_b$, whose support consists of graphs isomorphic to $G_b$ (that differ only in vertex and edge labelings), and each distribution is uniform over its support.
% \commentt{Do we think of $G_b$ as unlabeled as opposed to the graphs in the support of the distributions or is it also labeled? Or is $G_b$ a graph with specific labelings, i.e., one in the support? it is just a formality}
% We don't have time to get into such formalities..
We shall show that for a sufficiently large constant $c$, an algorithm that performs less than $n^{1/4}/c$ queries, cannot distinguish with success probability at least $2/3$ between a graph selected according to $\calD_0$ and a graph selected according to $\calD_2$. The lower bound stated in the theorem follows.

We start by defining $G_0$. The vertices in $G_0$ are partitioned into three sets: $X_0,Y_0,Z_0$, \SodaReb{whose sizes are $x_0,y_0$ and $z_0$, respectively. These sizes satisfy the following. First,}
 $x_0 = y_0(y_0-1)/2$. We set $y_0$ to be the largest odd integer satisfying 
 \ICALPCR{$x_0 + y_0 \leq n$}
 % $x_0+2y_0 \leq n$ (the reason for the factor 2 will become clear shortly), 
 and $z_0 = n-(x_0+y_0)$.
For each pair of vertices in $Y_0$ there is a unique vertex in $X_0$ that neighbors exactly these two vertices.
% Each vertex in $X_1$ has exactly two neighbors in $Y_1$, where any two vertices in $X_1$ share at most one neighbor in $Y_1$. Similarly, each vertex in $X_2$ has exactly two neighbors in $Y_2$, where any two vertices in $X_2$ share at most one neighbor in $Y_2$.
The vertices in $Z_0$ have no incident edges.
% Hence, all vertices in $Z_0$ have degree 0, all those in $X_0$ have degree 2, and all
% those in $Y_0 $ has degree $y_0 - 1$.
By construction, $G_0$ is $C_4$-free, and we have $y_0 = \Theta(n^{1/2})$, 
\ICALPCR{$z_0 = O(n^{1/2})$,}
%$z_0 = \Theta(y_0)=\Theta(n^{1/2})$, 
and $x_0 = \Theta(n)$.

Turning to $G_1$, its vertices are also partitioned into three sets: $X_1,Y_1,Z_1$, of sizes $x_1$, $y_1$ and $z_1$, respectively, where $X_1$ is further partitioned into two equal-size sets $X_1^1$ and $X_1^2$, and $Y_1$ is partitioned into two equal-size sets $Y_1^1$ and $Y_1^2$ (so that $x_1$ and $y_1$ are both even).
Here $x_1/2 = (y_1/2)^2$ so that $x_1 = y_1^2/2$, and we set 
% $y_1 = 2(y_0-1)$ 
\ICALPCR{
$y_1 = y_0 - 1$
}
and $z_1 = n- (x_1+y_1)$. For each pair of vertices, $u^1 \in Y_1^1$ and $u^2 \in Y_1^2$ there is one vertex in $X_1^1$ that neighbors exactly these two vertices, and one vertex in $X_1^2$ that neighbors them. The vertices in $Z_1$ have
no incident edges. By construction, since there are $x_1/2 = \Theta(n)$ edge-disjoint $C_4$s in $G_1$, it is $\Omega(1)$-far from being $C_4$-free.

Comparing the two graphs we observe that the vertices in
$Z_1$ have degree 0, which is the same as those in $Z_0$, The vertices in $X_1$ have degree 2, which is the same as those in $X_0$, and  the vertices in $Y_1$ have degree 
%$y_1/2 = y_0-1$, 
\ICALPCR{$2\cdot (y_1/2) = y_1 = y_0-1$,}
which is the same as those in $Y_0$.
Comparing the sizes of the sets of vertices, since 
%$y_1 = 2(y_0-1)$, 
\ICALPCR{
$y_1 = y_0 - 1$,
}
we get that
$$x_1 = \frac{y_1^2}{2} = \frac{(y_0-1)^2}{2} = \frac{y_0(y_0-1)}{2} - \frac{y_0-1}{2} = x_0 - \Theta(n^{1/2})\;,$$
and
% $$z_1 = n-(x_1+y_1) = n - (x_0 + (3/2)(y_0-1)) = z_0 - (y_0/2 - 3/2)) = \Theta(n^{1/2})$$
\ICALPCR{
$$z_1 = n-(x_1+y_1) = n - (x_0 + (1/2)(y_0-1)) = z_0 + y_0/2 + 1/2 = O(n^{1/2})\;.$$
}
%(recall that $x_0 + 2y_0 \leq n$). 
\ForFuture{D: maybe explain a bit more}
% Hence, $z_0 = O(n^{1/2})$, $y_0 = \theta(n^{1/2})$ and $x_0 = \Theta(n)$. Each vertex in $X_1$ has exactly two neighbors in $Y_1$, where any two vertices in $X_1$ share at most one neighbor in $Y_1$. Similarly, each vertex in $X_2$ has exactly two neighbors in $Y_2$, where any two vertices in $X_2$ share at most one neighbor in $Y_2$. The vertices in $Z_0$ an no incident edges. By construction, $G_0$ is $C_4$-free.

We next define two processes, $\calP_0$ and $\calP_1$, where $\calP_b$, for $b\in \bitset$ constructs ``on the fly'' a graph selected  according to $\calD_b$ while answering the queries of a testing algorithm. We shall show that unless one of several events occurs, the distributions induced on the sequences of queries and answers are identical. Let $q^t$ denote the $t$th query of the algorithm and $a^t$ the answer it gets, so that both are random variables. The query $q^t$ may depend on the \emph{query-answer history} $(q^1,a^1),\dots,(q^{t-1},a^{t-1})$ and possibly randomness of the algorithm. Both processes maintain a \emph{Knowledge graph}, which contains the vertices and edges that appear in these questions and answers, where $K^t=(V^t,E^t)$ is the knowledge graph after the $t$th query is answered. Let $X^t$, $Y^t$ and $Z^t$ denote the subsets of vertices with degree $2$, $y_0-1$ and $0$, respectively.
(In principle, the knowledge graph should also include information about the degrees of vertices and pairs of vertices that are known not to have edges between them, but, as we discuss below, this will be implicit from the graph.)
To determine $a^t$ given $q^t$,
the process $\calP_b$ simply considers all graphs in the support of $\calD_b$ that are consistent with $K^{t-1}$, and selects an answer with probability proportional to the number of such graphs that are also consistent with this answer.

We assume that before asking a query concerning a vertex $v$ that is not in the knowledge graph, the algorithm performs a degree query $q^t$ on $v$ (which in particular determines whether $v$ belongs to $X^t$, $Y^t$ or $Z^t$).
If the answer to this query is $2$, then the algorithm performs two neighbor queries to get $v$s two neighbors.
Also, if the algorithm gets a new vertex $v$ as an answer to a neighbor query from a vertex
$u \in Y^{t-1}$, then it performs a neighbor query on $v$ so as to gets its second neighbor.
A lower bound of $T$ on the query complexity of the algorithm under theses restrictions translates into a lower bound of $\Omega(T)$ without these restrictions. Note that the algorithm need not ask any further degree queries, as their answers are determined by the knowledge graph (which is known to be isomorphic to a subgraph of $G_0$ or $G_1$). Also note that because for each degree-2 vertex that is introduced into the graph, the algorithm obtains both its neighbors, we may assume (without loss of generality) that the algorithm does not perform any vertex-pair queries (as their answers would always be negative).


We observe that for both processes, if the algorithm performs $o(n^{1/2})$ queries, then the probability that any answer to a degree query (on a new vertex) is not $2$, is $o(1)$ (since $|Y_b|+|Z_b| = O(n^{1/2})$ for both $b=0$ and $b=1$). Let's refer to this ``bad'' event as $E_1$.
% From this point on we condition on $E_1$ not occurring. That is, all new vertices belong to $X^t$.
In addition, by the second restriction on the algorithm (i.e., that each degree-2 vertex essentially ``arrives with both its neighbors''), we have that for any neighbor query $q^t$ on a vertex in $Y^{t-1}$, the process $\calP_b$ answers with a vertex (in $X^t$) whose label is uniformly selected among all $n-|V^{t-1}|$ ``free'' labels. That is, the distribution on these answers is identical under both processes.


% We shall also assume that once a new vertex of degree $2$ (i.e., that belongs to $X^t$) is introduced into the knowledge graph (either following a degree query or an answer to a neighbor query for a vertex $u \in Y^{t-1}$), both its neighbors (in $Y^t$) are determined and provided to the algorithm ``for free''. Clearly, a lower bound  under this assumption provides a lower bound that is at least as high without this assumption.
% This assumption implies that there are no edges between vertices in the knowledge graph other than those that already appear in it.
% Hence,


It remains to analyze the answers to neighbor queries for vertices $v$ in $X^{t-1}$.
% choice of  neighbors of a newly introduced vertex $v$ in $X^t$.
If $v$ is the answer to a neighbor query from a vertex $u \in Y^{t-1}$, then each of the processes needs to select a single (additional) neighbor. Conditioned on this neighbor not belonging to $V^{t-1}$, for both processes, its label is uniformly distributed among all $n-|V^{t-1}|$ ``free'' labels. The same is true if $v$ is introduced following a degree query and both its neighbors are determined not to belong to $V^{t-1}$. We denote the event that we get at least one neighbor that already belongs to the knowledge graph, by $E_2$. Observe that conditioned on $E_1$ and $E_2$ not occurring,  the distributions on query-answer histories induced by the two processes, are identical (since, conditioned on these events, every degree query is answered by $2$, and every neighbor query is answered by a uniformly selected new label).

It hence remains to upper bound the probability of the event $E_2$, under both processes, when the algorithm performs at most $n^{1/4}/c$ queries, for some sufficiently large constant $c$. Consider any prefix $(q^1,a^1),\dots,(q^{t-1},a^{t-1})$ of the query-answer history (where neither event $E_1$ or $E_2$ has yet occurred), and the corresponding knowledge graph $K^{t-1}$.
Suppose that $q^t$ is a degree query for a vertex $v$ (that does not belong to $V^{t-1}$, and such that the answer to this query is $2$). Consider first the process $\calP_0$. 
 Since for each pair of vertices in $Y_0$ there is a single vertex in $X_0$ that neighbors them, and there are at most $2(t-1)$ vertices in $Y^{t-1}$ \ICALPCR{and at most $(t-1)$ in $X^{t-1}$}, the probability (over the distribution induced by $\calP_0$) that at least one of the neighbors of $v$ belongs to $Y^{t-1}$ is upper bounded by
\ICALPCR{
$$\frac{2(t-1)\cdot (y_0-1)}{x_0 - (t-1)} = \frac{2(t-1)\cdot (y_0-1)}{(y_0\cdot(y_0-1)/2) - (t-1)} %< \frac{4t}{y_0} 
    \leq \frac{c' t}{n^{1/2}}$$
    }
 for a constant $c'$ (since $y_0 = \Theta(n^{1/2})$ and $t = O(n^{1/4})$). A similar (slightly lower) upper bound is obtained for $\calP_0$ when $v$ is the answer to a neighbor query from a vertex $u \in Y^{t-1}$ (since one of $v$'s neighbors is already determined to be $u$).

Turning to $\calP_1$, we get a similar upper bound. Namely, for a new vertex $v$ (with degree 2), the probability that at least one of its neighbors belongs to $Y^{t-1}$ is upper bounded by
\ICALPCR{
$$\frac{2(t-1)\cdot (y_1/2)}{x_1-(t-1)} =\frac{2(t-1)\cdot (y_1/2)}{(y_1/2)^2-(t-1)} 
% = \frac{8(t-1)}{y_1} <\frac{4t}{y_0}
%= \frac{4(t-1)}{y_1} < \frac{4t}{y_0}
\leq \frac{c'' t}{n^{1/2}}\;,$$
}
for a constant $c''$,
 and a similar upper bound holds when $v$ is an answer to a neighbor query from $u\in Y^{t-1}$.
% (The reason we got the same final upper bound for both $\calP_0$ and $\calP_1$ despite the distributions being different is that in both cases we over-counted the number of possible pairs of neighbors that intersect $Y^{t-1}$ for a new degree-2 vertex.)
  Summing over $t=1,\dots,T= n^{1/4}/c$, we get that the probability that $E_2$ occurs (for either processes) is a small constant, as required.
\end{proof}

\subsection{A one-sided error lower bound for graphs with arboricity $\alpha$}
We now turn to a lower bound on non-constant arboricity graphs for one-sided error algorithms \ICALPCR{as stated in Theorem~\ref{thm:C4-gen-alpha-lb} in the introduction}. We assume that $\alpha =O(n^{1/2})$ (recall that every graph with  arboricity greater than $n^{1/2}$ must contain $C_4$s). %\commentt{DR: The lower bounds hold when $\alpha$ is known, and for the upper bounds we'll address the case of not knowing $\alpha$. Do we need to discuss the constant-factor gap between $n^{1/2}/c'$ and $n^{1/2}$?}  
The lower bound holds for $\alpha = \Omega(\log n)$. Note that the lower bound of Theorem~\ref{thm:C4-const-alpha-lb} also holds for $\alpha = O(\log n)$ that is non-constant, since, as noted in the introduction  we can add to the graphs in the lower-bound construction a subgraph with arboricity $\alpha$ that is $C_4$-free. \ForFuture{D: reference to Fox? or something else}

In order to prove Theorem~\ref{thm:C4-gen-alpha-lb}, we define
a distribution %$\calD$
over $n$-vertex graphs for which the following holds. First, with high probability, a graph selected according to %$\calD$
this distribution is $\Omega(1)$-far from being $C_4$-free. Second, the number of queries necessary to observe a $C_4$ with sufficiently high constant probability when the queries are answered by graph selected according to
the distribution %$\calD$
is $\Omega(n^{1/4}\alpha^{1/2})$.
Furthermore, if a one-sided error algorithm performs less than 
$n^{1/4}\alpha^{1/2}/c$ queries for a sufficiently large $c$ and does not observe a $C_4$, then it must accept, since the subgraph observed (including the degrees of the vertices) can be extended to a $C_4$-free graph.
% We start by describing $\calD$.

%%%%%%%%%%%%%%%%%
\iffalse
We shall define the distribution in two steps. First we define a distribution $\calD$ over graphs with multiple (parallel) edges, and then we use it to define a distribution $\calD'$ over graphs without multiple edges. \commentt{DR: should we just use the term parallel edges? or is it Ok that we use both terms?}
\fi
%%%%%%%%%%%%%%

%%%%%%%%%%%%%%%%%%%%%%%%%%%%%%%%
\iffalse
The vertex set of each graph in the support of $\calD$ is partitioned into five parts: $X_1$, $X_2$, $Y_1$, $Y_2$, and $Z$, where $|X_1| = |X_2|=n/4$, $|Y_1| = |Y_2| = n^{1/2}\alpha/4$ and $|Z| = n/2- n^{1/2}\alpha/2$.\footnote{We assume for simplicity that $n$ is divisible by 4 and that $n^{1/2}$ is an integer -- the more general case can be handled similarly by rounding.}
There is a bipartite graph between $X_b$ and $Y_{b'}$ for $b,b' \in \{1,2\}$, where vertices in $X_b$ have degree $\alpha$, vertices in $Y_{b'}$ have degree $n^{1/2}$ and vertices in $Z$ have degree $0$.  Each graph in the support of $\calD$ is defined by 4 matchings, each between two tables, where there are overall 8 tables: $T_{X_b,b'}$ for $b,b' \in \{1,2\}$ and $T_{Y_{b'},b}$ for $b,b' \in \{1,2\}$. For each for $b,b' \in \{1,2\}$, there is a matching between $T_{X_b,b'}$ and $T_{Y_{b'},b}$. The tables $T_{X_b,b'}$ are of dimensions $|X_b|\times \alpha$, and the dimensions of the tables $T_{Y_{b'},b}$, are of dimensions $|Y_{b'}|\times n^{1/2}$.

Denoting the vertices in $X_b$ by $x_{b,1},\dots,x_{b,|X_b|}$ and the vertices in $Y_{b'}$ by $y_{b',1},\dots,y_{b',|Y_b'|}$,
if entry $(i,j)$ in the first table is matched to entry $(i',j')$ in the second table, then
$T_{X_b,b'}[i,j] = (i',j')$ and $T_{Y_{b'},b}[i',j'] = (i,j)$, where
this means that there is an edge between $x_{b,i}\in X$ and $y_{b',i'}\in Y$, where this edge is the $((b-1)\cdot\alpha+j$)th edge incident to $x_{b,i}$ and the $(b'-1)\cdot n^{1/2}+j'$th edge incident to $y_{b',i'}$.
The distribution $\calD$ is uniform over the graphs in its support.

We define the \emph{extra multiplicity} of a graph $G$ as the total number of edges in $G$, minus the number of distinct edges. Thus, if between two vertices there are $\ell$ parallel edges, then this pair of vertices contributes $\ell-1$ to the extra multiplicity.
The distribution $\calD'$ (over graphs without multiple edges), is simply defined by the following process.
A graph $G$ is selected according to $\calD$, and it determines a graph $G' = \tau(G)$ by removing all occurrences of multiple edges but one per pair of vertices with at least one edge between them. The labels of the edges are modified so that they maintain the order of the first occurrence of each edge. \commentt{DR: Add example?}
Let $\Sup(\calD)$ ($\Sup(\calD')$) denote the support of $\calD$ ($\calD'$). \commentt{DR: some of what appears here might go to the preliminaries} Thus $\tau$ is a many-to-one and onto function from $\Sup(\calD)$ to $\Sup(\calD')$.
\fi
%%%%%%%%%%%%%%%%%%%%%%%%%%%%


The vertex set of each graph in the support of $\calD$ is partitioned into five parts: $X_1$, $X_2$, $Y_1$, $Y_2$, and $Z$, where $|X_1| = |X_2|=n/4$, and for $d = n^{1/2}/c_2$,  $|Y_1| = |Y_2| = (n\alpha)/(4 d) = \Theta(n^{1/2}\alpha)$ and $|Z| = n - (|X_1|+|X_2|+|Y_1|+|Y_2|)$. (We assume for simplicity that $n$ is divisible by $2d$
%4
and that $\alpha$ and $d$ are even integer -- the more general case can be handled similarly). % by rounding.) \commentt{DR: maybe need other similar assumptions, I think that $Y_b|$ should be divisible by $\alpha/2$ and  $|X_b|$ should be divisible by $d/2$}
% DR 5.6: I think the above comment is taken care of by requiring that $n$ be divisible by $2d$ where $d$ is even


The distribution $\calD$ is uniform over all graphs for which the following holds:
\begin{itemize}
\item There is a bipartite graph, denoted $G_{b,b'}$ between $X_b$ and $Y_{b'}$ for $b,b' \in \{1,2\}$, where vertices in $X_b$ have  $\alpha/2$ neighbors in $Y_{b'}$, vertices in $Y_{b'}$ have $d/2$ neighbors in $X_b$ and vertices in $Z$ have degree $0$.
\item There are no multiple edges.
\end{itemize}
 Observe that there exists at least one such graph: Let each $X_b$  consist of $\frac{n/4}{d/2}$ subsets, each of size $d/2$, and let each $Y_b$ consist of $\frac{n\alpha/(4d)}{\alpha/2}$ subsets, each of size $\alpha/2$. For the edges, let there be a complete bipartite graph between the $i$th subset in each $X_b$ to the $i$th subset of both $Y_1$ and $Y_2$. One can verify that indeed both conditions above hold. We stress that this graph is not a good candidate for a lower bound, as it contains many $C_4$s. However, as we prove below, when choosing uniformly over all graphs in the (non-empty) support of $\calD$, detecting a $C_4$ requires $\Omega(n^{1/4}\alpha^{1/2})$ queries.

% We shall  prove  that most graphs in the support of the distribution are $\Omega(1)$-far from being $C_4$-free, and then continue to prove that unless sufficiently many queries are performed, no cycle is detected with high probability. We start with the first claim.

 \begin{lemma}\label{lem:D-C4-far}
With probability at least $9/10$, a graph selected according to $\calD$ is $\Omega(1)$-far from being $C_4$-free.
\end{lemma}

 In order to establish Lemma~\ref{lem:D-C4-far},
 %To this end we shall
 we define the following ``bad events'' (taken over the random choice of a graph according to $\mathcal{D}$). Let $\alpha' = \Theta(\alpha)$ (where the constant in the $\Theta$ notation will be determined subsequently).
First, $E^0$ is the event that the fraction of edges between $X_1$ and $Y_1$ that participate in less than $\alpha'/2$ $C_4$s is larger than $1/2$.
% Conditioned on $E^0$ not occurring, there are initially at least  $|X_1|\cdot (\alpha/2)\cdot (\alpha'/2) = \Omega(n\alpha^2)$ $C_4$s.
Second, for each $b,b' \in \{1,2\}$ and $\ell \geq 1$, let $D_{b,b'}^\ell$ be the set (\emph{bucket}) of edges between $X_b$ and $Y_{b'}$ that participate in a number of $C_4$s that is in the range $[(2^\ell+1)\cdot \alpha', (2^{\ell+1}+1)\cdot \alpha')$.
The event $E_{b,b'}^\ell$ is the event that $|D_{b,b'}^\ell| > (2^{-2\ell}/c_3)\cdot |X_b|\cdot (\alpha/2)$ for some constant $c_3$.
% \commentt{DR: the reason for the annoying $+1$ in the definition of the buckets is because Chebishev bounds the additive deviation from the expected value, which is $\alpha'$. }

\begin{claim}~\label{clm:Pr-Es}
For a sufficiently large $\alpha = \Omega(\log n)$,
the probability that either $E^0$ occurs or some $E_{b,b'}^\ell$ occurs (for $b,b' \in \{1,2\}$ and $\ell \geq 1$) is at most $1/10$.
\end{claim}
% We shall show that (for sufficiently large $\alpha = \Omega(\log n)$), the probability of each of these $O(\log n)$ events is sufficiently small so that when we take a union bound over all of them we get that the probability that any one of them occurs is at most $1/10$.
We prove Claim~\ref{clm:Pr-Es} momentarily, and first show how it implies Lemma~\ref{lem:D-C4-far}.

\medskip
\begin{proofof}{Lemma~\ref{lem:D-C4-far}}
Assume that neither $E^0$ nor any $E_{b,b'}^\ell$ occur, which by Claim~\ref{clm:Pr-Es} holds with probability at least  $9/10$. Conditioned on $E^0$ not occurring, there are at least  $|X_1|\cdot (\alpha/2)\cdot (\alpha'/2)$ ($ = \Omega(n\alpha^2)$) $C_4$s in a graph selected according to $\mathcal{D}$.
%Conditioned on none of them occurring,
Suppose we mark the edges that participate in any one of the buckets $D_{b,b'}^\ell$ for $b,b'\in \{1,2\}$ and $\ell\geq 1$ by ``red'' and all other edges by ``green''.
By our assumption on the sizes of the buckets (the events $E_{b,b'}^\ell$ not occurring), the total number of $C_4$s that contain at least one red edge is upper bounded by
\ForFuture{ D: verify constants etc.}
\begin{equation}
\sum_{b,b'}\sum_\ell |D_{b,b'}^\ell|\cdot (2^{\ell+1}+1)\cdot \alpha' \leq
  4 \cdot (2^{-2\ell}/c_3)\cdot |X_b|\cdot (\alpha/2) \cdot (2^{\ell+1} +1)\cdot \alpha' \leq   |X_1|\cdot(\alpha/2)\cdot (\alpha'/c_3') %= O(n\alpha^2)
  \;.
\end{equation}
By ensuring the $c_3' \geq 4$,
% conditioned on all bad events not occurring,
the number of $C_4$s residing  on green edges (only) is at least
\begin{equation}
\frac{|X_1|\cdot\alpha\cdot\alpha'}{4}
- \frac{|X_1|\cdot\alpha\cdot\alpha'}{8} =
\Omega(n \alpha^2)\;.
\end{equation}
%, as desired.

Hence, with probability at least $9/10$ over the choice of a graph according to $\calD$, the following holds. There exists a set of green edges,  such that the number of $C_4$s that reside  on green edges is $\Omega(n\alpha^2)$, and each green edge participates in  $O(\alpha)$ $C_4$s. This implies that there are $\Omega(n\alpha)$ edge-disjoint $C_4$s. To verify this consider running an iterative procedure, starting from the graph restricted to green edges. In each iteration, until there are no remaining $C_4$s, a $C_4$ is selected from the current graph and put in the set of edge-disjoint $C_4$s. Then all $O(\alpha)$ $C_4$s that include one of the edges in the selected $C_4$ are removed from the graph (i.e., their other edges are removed from the graph). Since we start with $\Omega(n\alpha^2)$ $C_4$s and in each iteration remove $O(\alpha)$ $C_4$s, the final set of edge-disjoint $C_4$s we get is of size $\Omega(n\alpha)$, implying that the graph is $\Omega(1)$-far from $C_4$-free.
\end{proofof}

\smallskip
\noindent
 We next prove Claim~\ref{clm:Pr-Es}.

 \smallskip
\begin{proofof}{Claim~\ref{clm:Pr-Es}}
We start by noting that in the selection of a random graph according to $\mathcal{D}$, the four graphs $G_{b,b'}$, $b,b' \in \{1,2\}$ are generated independently from each other (recall that $G_{b,b'}$ is a bipartite graph between $X_b$ and $Y_{b'}$, where each vertex in $X_b$ has $\alpha/2$ neighbors in $Y_{b'}$, and each vertex in $Y_{b'}$ has $d/2$ neighbors in $X_b$).

\medskip
For simplicity of notations, let $b=1$, $b'=1$.
The analysis below can be applied to any pair $b,b' \in \{1,2\}$.
Fix any choice of bipartite graph $G_{2,2}$.
%between $X_2$ and $Y_2$ such that every vertex in $X_2$ has $\alpha/2$ neighbors in $Y_2$ and every vertex in $Y_2$ has $d/2 = n^{1/2}/(2c_2)$ neighbors in $X_2$.
We shall show that for any fixed edge $(x,y)$ between $x\in X_1$ and $y \in Y_1$, the following holds over the random choice of $G_{1,2}$ and $G_{2,1}$.
%when we consider the random choice of the bipartite graph between $X_1$ and $Y_2$ and the bipartite graph between $Y_1$ and $X_2$.
The expected number of $C_4$s that reside on $(x,y)$ is $\alpha'$, and furthermore, with sufficiently high probability, the actual number does not deviate by too much from its expectation.

For a fixed edge $(x,y)\in X_1\times Y_1$,
% consider the set of $d/2$ neighbors of $y$ in $X_2$,
Let $\Gamma_2(x)$ denote  the set of $\alpha/2$ neighbors of $x$ in $Y_2$ (as determined by the choice of $G_{1,2}$), and let $\Gamma_2(y)$ denote the set of $d/2$ neighbors of $y$ in $X_2$ (as determined by the choice of $G_{2,1}$). Since we are considering a single vertex $x\in X_1$ and a single vertex $y\in Y_1$, these sets are uniformly distributed in $Y_2$ and $X_1$, respectively.

For $i\in [d/2]$  and $j \in [\alpha/2]$, let $\chi_{i,j}$ be a Bernoulli random variable that is 1 if there is an edge between the $i$th neighbor of $y$ and the $j$th neighbor of $x$ (and is 0 otherwise). Let  $\mu = \Pr[\chi_{i,j}]=1]$, so that
\begin{equation}
\mu = % \frac{|X_2|\alpha}{|X_2||Y_2|} =
   \frac{\alpha/2}{|Y_2|} %= \frac{d/2}{|X_2|}
   = \Theta\left(\frac{n^{1/2}}{n}\right) = \Theta\left(\frac{1}{n^{1/2}}\right)\;.
\end{equation}
Therefore, if we let $\alpha' = (d/2)\cdot(\alpha/2)\cdot \mu$ (so that $\alpha' = \Theta(\alpha)$), then
\begin{equation}
\Ex\left[\sum_{\substack{i\in [d/2] \\ j\in [\alpha/2]}}\chi_{i,j}\right] = \frac{d}{2}\cdot\frac{\alpha}{2}\cdot \mu = \alpha'\;.
\end{equation}
We turn to analyze the variance of the sum.
\begin{align}
\Var\left[\sum_{\substack{i\in [d/2] \\ j\in [\alpha/2]}}\chi_{i,j}\right] & =
\Ex\left[\left(\sum_{\substack{i\in [d/2] \\ j\in [\alpha/2]}}\chi_{i,j}\right)^2\right] - \left(\Ex\left[\sum_{\substack{i\in [d/2] \\ j\in [\alpha/2]}]}\chi_{i,j}\right]\right)^2 \\
& =
\Ex\left[\sum_{\substack{i\in [d/2] \\ j\in [\alpha/2]}} (\chi_{i,j})^2\right] +
\Ex\left[\sum_{\substack{i\in [d/2] \\ j_1\neq j_2 \in [\alpha/2]}}\chi_{i,j_1}\chi_{i,j_2}\right]
% \nonumber\\
%&\;\;\;\;\;\;\;\;\;\;
+ \Ex\left[\sum_{\substack{i_1\neq i_2\in [d/2] \\ j\in  [\alpha/2]}}\chi_{i_1,j}\chi_{i_2,j}\right] \nonumber \\
&\;\;\;\;\;\;\;\;\;\;+ %2
\Ex\left[\sum_{\substack{i_1\neq i_2\in [d/2] \\ j_1\neq j_2\in  [\alpha/2]}}\chi_{i_1,j_1}\chi_{i_2,j_2}\right] - \left(\frac{d}{2}\cdot\frac{\alpha}{2}\cdot \mu\right)^2 \;.
\label{eq:var-lb}
\end{align}
% \commentt{DR: there was a factor of 2 before the last expectation and I don't think there should be - the expression is already counting, say both $i_1=1$, $i_2=2$ and $i_1=2$, $i_2-1$}
We analyze each of the terms separately.
Since the $\chi_{i,j}$s are Bernoulli random variables,
$(\chi_{i,j})^2 = \chi_{i,j}$ so that
\begin{equation} \label{eq:cheb-C4-lb-case1}
\Pr[(\chi_{i,j})^2 = 1] = \Pr[\chi_{i,j}=1] = \mu\;.
% \sum_{\substack{i \in [d/2]\\ j\in [\alpha/2]}} (\chi_{i,j})^2 = \sum_{i \in [d/2], j\in [\alpha/2]} \chi_{i,j} =
% \frac{d}{2}\cdot\frac{\alpha}{2}\cdot \mu\;.
\end{equation}
% \commentt{DR: should probably continue with $\alpha'$. Moreover, maybe we can give a more general Chebishev based statement that can also be applied here}
 For every $i\in [d/2]$ and $j_1,j_2 \in [\alpha/2]$ such that $j_1 \neq j_2$,
 % since  $\Pr[\chi_{i,j_2}=1 \,|\, \chi_{i,j_1}=1]
  %\leq \Pr[\chi_{i,j_2}=1]$,
\begin{equation}\label{eq:cheb-C4-lb-case2}
\Pr[\chi_{i,j_1}=\chi_{i,j_2}=1]
  \leq \left( \frac{\alpha/2}{|Y_2|}\right)^2 = \mu^2
  %= O(1/n)
  \;,
\end{equation}
and for every $i_1,i_2\in [d/2]$ and $j\in [\alpha/2]$ such that $i_1\neq i_2$,
\begin{equation}\label{eq:cheb-C4-lb-case3}
 \Pr[\chi_{i_1,j}=\chi_{i_2,j}=1]
\leq %\Ex[\chi_{i_1,j_1}\cdot \chi_{i_2,j_1}] = \Pr[\chi_{i_1,j_1}=\chi_{i_1,j_2}]
\left( \frac{d/2}{|X_2|}\right)^2 = \left( \frac{\alpha/2}{|Y_2|}\right)^2 = \mu^2 %= O(1/n)
\;.
\end{equation}
Finally,
for $i_1\neq i_2\in [d/2]$ and $j_1\neq j_2\in [\alpha/2]$, we have that
\begin{equation}\label{eq:cheb-C4-lb-case4}
\Pr[\chi_{i_1,j_1}=\chi_{i_2,j_2}=1]
=  \Pr[\chi_{i_1,j_1}=1] \cdot
\Pr[\chi_{i_1,j_1}=1\,|\,\chi_{i_2,j_1}=1]
  =\frac{\alpha/2}{|Y_2|}\cdot \frac{|X_2|(\alpha/2)-1}{(|X_2|-1)(|Y_2|-1)}\;,
\end{equation}
where the second term on the right-hand-side of Equation~\eqref{eq:cheb-C4-lb-case4} is the fraction of pairs in $X_2\times Y_2$ that have an edge between them, among all pairs with the exception of one pair that is an edge in $G_{2,2}$.
%probability  of obtaining an edge in $G_{2,2}$, when selecting a vertex from $X_2$ and a vertex from $Y_2$
It is not hard to verify that
\begin{equation}\label{eq:cheb-C4-lb-5}
 \frac{|X_2|(\alpha/2)-1}{(|X_2|-1)(|Y_2|-1)} \leq \frac{\alpha/2}{|Y_2|} \cdot \left(1+\frac{2}{|Y_2|}\right) = \mu \cdot \left(1+\frac{2}{|Y_2|}\right)
\end{equation}
By combining Equations~\eqref{eq:var-lb}--\eqref{eq:cheb-C4-lb-5} (and recalling that $\alpha' = (d/2)(\alpha/2)\mu$), we get that
\begin{eqnarray}
\Var\left[\sum_{\substack{i\in [d/2] \\ j\in [\alpha/2]}}\chi_{i,j}\right]
&\leq& \frac{d}{2}\cdot\frac{\alpha}{2}\cdot\mu
+ \frac{d}{2}\cdot \left(\frac{\alpha}{2}\right)^2 \cdot \mu^2 + \left(\frac{d}{2}\right)^2\cdot \frac{\alpha}{2}\cdot \mu^2 \nonumber \\
&\;\;\;\;\;\;\;\;\;\;\;\;\;\;\;\;\;\;\;+& \left(\frac{d}{2}\right)^2\cdot \left(\frac{\alpha}{2}\right)^2 \cdot \mu^2\cdot \left(1+\frac{1}{|Y_2}\right) - \left(\frac{d}{2}\cdot\frac{\alpha}{2}\cdot \mu\right)^2 \nonumber \\
&=& \alpha'\cdot
\left( 1+ \frac{\alpha}{2}\cdot\mu + \frac{d}{2}\cdot \mu + \frac{d}{2}\cdot \frac{\alpha}{2}\cdot \frac{2}{|Y_2|}\cdot \mu\right)\nonumber \\
&\leq& 4\alpha' \;,
\end{eqnarray}
where in the last inequality we used the following facts: $\alpha \leq d$, $|Y_2|> \alpha$, and
$(d/2)\mu \leq 1$.

By Chebishev's inequality,
%Claim~\ref{clm:var1},
for any choice of $t>0$,
%(where we shall consider $t=2^{\ell}$ for $\ell = -1,0,1,2\dots$),
\begin{eqnarray}
\Pr\left[\left|\sum_{i\in [d/2],j\in [\alpha/2]}\chi_{i,j} - \alpha'\right|< t\cdot \alpha' \right] &\leq& \frac{\Var\left[\sum_{\i\in [d/2],j\in [\alpha/2]}\chi_{i,j}\right]}{t^2\cdot \left(\Ex\left[\sum_{i\in [d],j\in [\alpha]} \chi_{i,j}\right]\right)^2} \\
&\leq& \frac{4\alpha'}{t^2 \cdot (\alpha')^2}
% \\
%&\leq& \frac{\alpha' \cdot (1+(\alpha/2)\mu + (d/2)\mu)}{t^2 \cdot (\alpha')^2} \; < \; \frac{3}{t^2 \alpha'}\;,
\;\;= \;\;\frac{4}{t^2\cdot \alpha'}\;
\label{eq:C4-lb-cheb2}
\end{eqnarray}
By setting $t=1/2$ we get that for any fixed edge $(x,y)\in X_1\times Y_1$, the probability that $(x,y)$ participates in less than $\alpha'/2$ $C_4$s is at most $16/\alpha'$. By Markov's inequality and the definition of the event $E^0$, we have that
$\Pr[E^0] \leq 32/\alpha'$.
% \commentt{for some reason was 100 before, not really important of course, but should doublecheck}

By applying the same argument to any other $b,b' \in \{1,2\}$ (i.e., not necessary $b=b'=1$), we get that
for any fixed edge $(x,y)\in X_b\times Y_{b'}$ and for
any choice of $t = 2^{\ell}$, $\ell \geq 0$,
 the probability that $(x,y)$ participates in more  than $(2^\ell +1)\cdot \alpha'$ $C_4$s is at most $\frac{4}{2^{2\ell}\cdot \alpha'}$.
Therefore, the expected number of such edges is at most
$|X_b|\cdot (\alpha/2)\cdot \frac{4}{2^{2\ell}\cdot \alpha'}$. By Markov's inequality, the probability that there are more than
$ |X_b|\cdot (\alpha/2)\cdot \frac{1}{c_3\cdot 2^{2\ell}}$ such edges is at most $\frac{4c_3}{\alpha'}$.
% $D_{b,b'}^\ell$ be the set (\emph{bucket}) of edges between $X_b$ and $Y_{b'}$ that participate in a number of $C_4$s that is in the range $[(2^\ell+1)\cdot \alpha', (2^{\ell+1}+1)\cdot \alpha')$.
%The event $E_{b,b'}^\ell$ is the event that $|D_{b,b'}^\ell| > (2^{-2\ell}/c_3)\cdot |X_b|\cdot (\alpha/2)$ for some constant $c_3$.

By the definition of the bucket $D_{b,b'}^\ell$ (which contains a subset of the aforementioned edges) and the event $E^\ell_{b,b'}$, we have that $\Pr[E_{b,b'}^\ell] \leq \frac{4c_3}{\alpha'}$.
Since $\alpha' = \Theta(\alpha)$ and $\alpha \geq c_1\log n$, for a sufficiently large constant $c_1$, if we take a union bound over all bad events, we get a total probability of at most $1/10$, and the claim follows.
\end{proofof}
%%
%%%%%%%%%
\iffalse
Let $\bar{\chi}_{i,j} = \chi_{i,j} - \mu$, so that $\Ex[\bar{\chi}_{i,j}] = 0$.
By Chebishev's inequality, for any choice of $t$ (where we shall consider $t=1/2,2,4\dots$),
\begin{equation}\label{eq:cheb}
\Pr\left[\left|\sum_{i\in [d],j\in [\alpha]}\chi_{i,j} - d\cdot \alpha\cdot \mu\right|< t\cdot d\cdot \alpha\cdot \mu \right] \leq \frac{\Var\left[\sum_{i\in [d],j\in [\alpha]}\chi_{i,j}\right]}{t^2\cdot \left(\Ex\left[\sum_{i\in [d],j\in [\alpha]} \chi_{i,j}\right]\right)^2}
= \frac{\sum_{i_1,j_2,i_2,j_2} \Ex[\bar{\chi}_{i_1,j_1}\cdot \bar{\chi}_{i_2,j_2}]}{t^2 \alpha^2 d^2 \mu^2} \;.
\end{equation}
Since $d \cdot \mu = \Theta(1)$, the  denominator is $\Theta(t^2 \alpha^2)$.
% we detect a triangle when selecting $\theta$ random neighbors of $a$ and $\theta$ random neighbors of $b$.

In order to bound the numerator,
we break the sum $\sum_{i_1,j_2,i_2,j_2} \Ex[\bar{\chi}_{i_1,j_1}\cdot \bar{\chi}_{i_2,j_2}]$ into four sub-sums.
\paragraph{The first} is over $i_1\neq i_2$ and $j_1 \neq j_2$. For such quadruples of indices, $\chi_{i_1,j_1}$ is independent of $\chi_{i_2,j_2}$ and so
\begin{equation}\label{eq:cheb-C4-lb-case1}
\Ex[\bar{\chi}_{i_1,j_1}\cdot \bar{\chi}_{i_2,j_2}] =  \Ex[\bar{\chi}_{i_1,j_1}]\cdot \Ex[\bar{\chi}_{i_2,j_2}] = 0\;.
\end{equation}
% where the last iequality follows from the definition of $\bar{\chi}_{i,j}$.
Hence, the first sub-sum does not contribute anything to the right-hand-side of Equation~\eqref{eq:cheb}.

\paragraph{The second} is over $i_1 = i_2$ and $j_1 = j_2$, so that
\begin{equation}\label{eq:cheb-C4-lb2-case2}
\Ex[\bar{\chi}_{i_1,j_1}\cdot \bar{\chi}_{i_2,j_2}] = \Ex[(\bar{\chi}_{i_1,j_1})^2]
\leq \Ex[(\chi_{i_1,j_1})^2] = \Ex[\chi_{i_1,j_1}] = \mu \;,
\end{equation}
where the second inequality follows from the definition of $\bar{\chi}_{i,j}$.
Since there are $d\cdot \alpha$ quadruples in this sub-sum, it contributes $d\cdot\alpha\cdot \mu = O(\alpha)$.
% after dividing by the denominator of the right-hand-side of
% Equation~\eqref{eq:cheb}, we get $\frac{\theta^2\mu}{\theta^4\mu^2} = \frac{1}{\theta^2 \mu}$.
% By Equation~\eqref{eq:mu-lb} and the setting of $\theta$, this sub-sum contributes at most $1/6$.

\paragraph{The third} is over $i_1 = i_2$ and $j_1 \neq j_2$, in which case we have
\begin{equation}\label{eq:cheb-case2}
%\Ex[\bar{\chi}_{i_1,j_1}\cdot \bar{\chi}_{i_2,j_2}]  =
\Ex[\bar{\chi}_{i_1,j_1}\cdot \bar{\chi}_{i_1,j_2}]
\leq \Ex[\chi_{i_1,j_1}\cdot \chi_{i_1,j_2}] = \Pr[\chi_{i_1,j_1}=\chi_{i_1,j_2}=1]
  \leq \left( \frac{\alpha}{|Y_2|}\right)^2 = \mu^2 = O(1/n)\;.
\end{equation}
\commentt{DR: double check this}
Since there are $d\cdot \alpha^2$ quadruples in this sub-sum, it contributes $O(\alpha \cdot (\alpha/n^{1/2}))$.

\paragraph{The fourth} is over $i_1 \neq i_2$ and $j_1 = j_2$, where similarly to the third case we get
\begin{equation}\label{eq:cheb-case3}
\Ex[\bar{\chi}_{i_1,j_1}\cdot \bar{\chi}_{i_2,j_2}]  %= \Ex[\bar{\chi}_{i_1,j_1}\cdot \bar{\chi}_{i_2,j_1}]
\leq %\Ex[\chi_{i_1,j_1}\cdot \chi_{i_2,j_1}] = \Pr[\chi_{i_1,j_1}=\chi_{i_1,j_2}]
\left( \frac{d}{|X_2|}\right)^2 = O(1/n)\;.
\end{equation}
Since there are $d^2\cdot \alpha = O(n\alpha)$ quadruples in this sub-sum, it contributes $O(\alpha)$.

Combining all sub-sums we get that the probability of deviating by more than $t$ times the expected value is upper bounded by $\frac{1}{t^2 \alpha}$.
\fi
%%%%%%%%%%%%%%%




% \commentt{DR: we might want to first introduce some smaller subclaims (maybe this should be a lemma and they will be claims)}


%
%%%%%%%%%%%%%%%%%%%%%%%%%%%
\iffalse
In order to prove the claim we first show that with probability at least $1-XX$ over the choice of a graph $G$ according to $\calD$, the graph $G$ contains at least $n\alpha/c$ \emph{edge-disjoint} $C_4$s (when counting multiple edges as distinct).
%  and it contains at most $n\alpha/(2c)$ multiple edges.
In addition we show that with probability at least $1-YY$ over the choice of a graph $G$ according to $\calD$,
% it contains
the extra multiplicity of $G$ (as defined above) is at most $n\alpha/(2c)$.
% multiple edges.
The claim follows by a union bound.

% We start by establishing the second subclaim. We shall actually prove a stronger statement: with probability at least $1-YY$, no edge has multiplicity of more than 2
% for every vertex $x \in X_2\cup X_2$, the number of distinct neighbors of $x$ in $Y_b$ for $b\in \{1,2\}$ is at least $\alpha/2$ and an analogous statement holds for  $y \in Y_1\cup Y_2$. To verify this consider any $x_{b,i}\in X_b$ and the matching of its corresponding entries $T_{X_b}[i,j]$ for $j=1,\dots,\alpha$.

% We now turn to the first subclaim, while we condition on the ``good events'' of the first subclaim holding for every $x\in X_1$ and $y\in Y_1$.
Consider first selecting the edges between $X_1$ and $Y_1$, between $X_1$ and $Y_2$ and between $X_2$ and $Y_1$.
%For each edge $(x,y) \in X_1\times Y_1$, where there are $|X_1| \cdot \alpha = \Theta(n\alpha)$ such edges, let
For each $x\in X_1$, let $\Gamma_2(x)$ denote the multiset of its neighbors in $Y_2$ (that is, if there are $t$ edges between $x$ and $y'$, then $y'$ appears $t$ times in $\Gamma_2(x)$), and similarly define $\Gamma_2(y)$ for each $y \in Y_2$. By construction, $|\Gamma_2(x)| = \alpha$ and $|\Gamma_2(y)| = n^{1/2}$.

Before continuing with adding edges between $X_2$ and $Y_2$, we show that with probability at least $1-XX'$, there is no edge in $(X_1\times Y_2) \cup (X_2 \times Y_1)$ whose multiplicity is larger than four. To this end, consider any vertex $x_{1,i} \in X_1$ (the same argument applies to $x_{2,i}\in X_2$). For any choice of five entries in row $i$ of $T_{X_1,2}$,
% quadruple of entries $T_{X_1}[i,j_1],T_{X_1}[i,j_2],T_{X_1}[i,j_3]$,
the probability that all these entries are matched to five entries
in the same row $i'$ of  $T_{Y_2,1}$, for some $i'$, is upper bounded by $(1/|Y_2|)^4 = O(1/(n^2\alpha^4)$.
The probability that this event occurs for any choice of five entries in row $i$ is hence $O(\alpha/n^2)$, and the probability that this occurs for some $i$ is $O(\alpha/n) = o(1)$.
From this point on we condition on the event, denoted $E_1$, that every edge in $X_1\times Y_2 \cup X_2 \times Y_1$ has multiplicity at most four.

Let $(x,y)$ be any fixed edge in $X_1\times Y_1$, where there are $|X_1| \cdot \alpha = \Theta(n\alpha)$ such edges.
Next consider the selection of edges incident to vertices in $\Gamma_2(y)$. Namely, assume without loss of generality that $\Gamma_2(y) = \{x_{2,1},\dots,x_{2,n^{1/2}}\}$, and $\Gamma_2(x) = \{y_{2,1},\dots,y_{2,\alpha}\}$. We match the entries $T_{X_2,2}[i,j]$ for $i=1,\dots,n^{1/2}$ and $j=1,\dots,\alpha$ to entries of $T_{Y_2,2}$ according to lexicographic order.
We define $s=|\Gamma_2(y)|\cdot \alpha = n^{1/2}\cdot \alpha$ random variables $\chi_1,\dots,\chi_s$, where $\chi_\ell = 1$ if the $\ell$th entry according to the above order is matched to an entry $T_{Y_2,2}[i',j']$ where $i'\in [\alpha]$ (so that $y_{2,i'}\in \Gamma_2(x)$). Let $\chi = \sum_{\ell=1}^s \chi_\ell$.  Observe that the number of $C_4$ that the edge $(x,y)$ participates in and contain a pair of vertices $x' \in X_2$ and $y' \in Y_2$, is lower bounded by $\chi$, and conditioned on the event $E_1$, it is upper bounded by
$16\cdot \chi$. \commentt{DR: maybe need to elaborate a bit more.}
While these random variables are not independent, we can bound the probability that $\chi \notin [\alpha/2,2\alpha]$
by considering related, independent random variables, as explained next.

%Specifically, %for the lower bound of $\alpha/2$,
Observe that for each $\ell \in [s]$
and for any $\beta \in \{0,1\}^{\ell-1}$,
$$\Pr[\chi_\ell=1 \,|\, \chi_1\dots,\chi_{\ell-1}=\beta] =
       \frac{\alpha\cdot n^{1/2} - \sum_{r=1}^{\ell-1}\beta_i}{\alpha\cdot n - (\ell-1))}\;.$$
Therefore, for any $\beta \in \{0,1\}^{\ell-1}$ such that $\sum_{r=1}^{\ell-1}\beta_i < \alpha/2$, we have that
$$\Pr[\chi_\ell=1 \,|\, \chi_1\dots,\chi_{\ell-1}=\beta] \geq
           \frac{\alpha\cdot n^{1/2} - \alpha/2} {\alpha\cdot n } = \frac{1}{n^{1/2}}-\frac{1}{2n}\;,$$
and for any $\beta \in \{0,1\}^{\ell-1}$
$$\Pr[\chi_\ell=1 \,|\, \chi_1\dots,\chi_{\ell-1}=\beta] \leq
    \frac{\alpha\cdot n^{1/2}}{\alpha\cdot n - \alpha\cdot n^{1/2}} = \frac{1}{n^{1/2}-1}\;.$$

We can hence obtain an upper bound on the probability that $\chi < \alpha/2$ by considering $s$ independent random variables $\chi'_1,\dots,\chi'_s$, where $\Pr[\chi'_\ell=1] = \frac{1}{n^{1/2}}-\frac{1}{2n}$ and applying the multiplicative Chernoff bound. We get that the probability that $\sum_{\ell=1}^s \chi'_\ell < \alpha/2$ is $\exp(-\Omega(\alpha))$, so that the same statement holds for $\chi$. For an upper bound on the probability that $\chi > 2 \alpha$, consider $s$ independent random variables $\chi''_2,\dots,\chi''_s$, where $\Pr[\chi''_\ell=1] = \frac{1}{n^{1/2}-1}$. Here too by applying the multiplicative Chernoff bound we get that the probability that $\sum_{\ell=1}^s \chi'_\ell > 2\alpha$ is $\exp(-\Omega(\alpha))$, and the same statement holds for $\chi$.

Since $\alpha = \Omega(\log n)$, by taking a union bound over all edges between $X_1$ and $Y_1$, we get that
with probability at least $1-XX$, every edge $(x,y)\in X_1\times Y_1$ participates in at least $\alpha/2$ $C_4$s with a vertex in $X_2$ and a vertex in $Y_2$, and at most $32\alpha$ such $C_4$s. We can hence get a set of at least $\alpha\cdot |X_1|/64 = \Omega(\alpha n)$ edges-disjoint $C_4$s by applying a simple iterative algorithm for constructing such a set. \commentt{DR: explain more/better?}

Finally, we  bound the extra multiplicity obtained.
% number of multiple edges \commentt{DR: need a better term: it is the number of times that we get an edge that we got before}.
The expected number of ``extra'' parallel edges is upper bounded by
$(n/2)\cdot {\alpha\choose 2} \cdot 4/ n^{1/2} \leq \alpha^2 n^{1/2}$. \commentt{TO BE finalized after setting parameters more precisely or maybe just upper bounding $\alpha$}
\fi
%%%%%%%%%%%%%%%%%%%%%%%%%
%\end{proof}
% setting of $\chi_1,\dots,\chi_{\ell-1} \in \{0,1\}^{\ell-1}$,
%%

\begin{lemma}\label{lem:D-C4-lb}
Let $\calA$ be an algorithm  that performs less than $n^{1/4}\alpha^{1/2}/c$ queries for a sufficiently large constant $c$. The probability that $\calA$ detects a $C_4$ in a graph selected according to $\calD$
% $\calD'$,
is at most $1/10$.
\end{lemma}

\begin{proof}
Similarly to the proof of Theorem~\ref{thm:C4-const-alpha-lb}, consider a process $\calP$ that answers the query of a testing algorithm while constructing a graph from $\Sup(\calD)$ (the support of $\calD$) on the fly. Here too we define a knowledge graph $K^t$ following the query-answer history $(q^1,a^1),\dots,(q^t,a^t)$ and we require the algorithm to perform a degree query on each vertex $v$ that is not yet in the knowledge graph before performing any other query involving it. The algorithm also receives the subset to which the new vertex belongs (as this is not be implied by the degree of the vertex). Let $X_1^t$, $X_2^t$, $Y_1^t$, $Y_2^t$ and $Z^t$ denote the corresponding subsets in the knowledge graph $K^t$. %\commentt{DR: assume $\alpha < n^{1/2}$?}

The set of edges in $K^t$ is denoted by $E^t$, where for each $e\in E^t$ there are two associated labels,
$\phi^t_X(e)$ for its endpoint in $X^t = X_1^t \cup X_2^t$, and $\phi^t_Y(e)$ for its endpoint in $Y^t = Y_1^t \cup Y_2^t$.
 We assume that a neighbor query $(x_{b,i},j)$ is answered by a pair $(y_{b',i'},j')$ where
$y_{b',i'}$ is the $j$th neighbor of $x_{b,i}$ and $x_{b,i}$ is the $j'$th neighbor of $y_{b',i'}$. Similarly, if a vertex-pair query $(x_{b,i},y_{b',i'})$ is answered positively, then the algorithm also gets the corresponding edge labels $\phi_X$ and $\phi_Y$ for this edge. The knowledge graph also includes a set $\overline{E}^t$ of non-edges (corresponding to negative answers to vertex-pair queries).

%%%%%%%%%%%
\iffalse
Recall that each graph $G' \in \Sup(\calD')$ may correspond to several graphs in $\Sup(\calD)$. Namely, to all graphs in $\tau^{-1}(G')$. Hence, the weight of $G'$ according to $\calD'$ is %the fraction of these graphs in the support of $\calD$.
$|\tau^{-1}(G')|/|\Sup(\calD)|$.
Therefore,
\fi
%%%%%%
Similarly to
the proof of Theorem~\ref{thm:C4-const-alpha-lb}, in order to
determine $a^t$ given $q^t$,
the process $\calP$ considers all graphs  $G\in \Sup(\calD)$
that are consistent with $K^{t-1}$,
%such that $\tau(G)$ is consistent with $K^{t-1}$.
It then selects an answer with probability proportional to the number of such graphs $G$ for which %$\tau(G)$
$G$ is also consistent with this answer.
% This is equivalent to selecting a graph $G$ in the support of $\calD$ uniformly among all those whose corresponding graph $G'$ in the support of $\calD'$ is consistent with $K^t$, and answering according to $G'$.

Consider any knowledge graph $K^{t-1}$ for $t < n^{1/4}\alpha^{1/2}/c$.
% so that $|X_1^{t-1}|$, $|X_2^{t-1}|$, $|Y_1^{t-1}|$, $|Y_2^{t-1}|$ $|E^{t-1}|$ and $|\overline{E}^{t-1}|$ are all smaller than $n^{1/4}\alpha^{1/2}/c$.
We would like to bound the probability that either a neighbor query or a vertex-pair query adds to the knowledge graph an edge $(x,y)$ that creates a $C_4$ in the knowledge graph. We shall actually bound the probability that a neighbor query is answered with \emph{some} vertex that already belongs to the knowledge graph or a vertex-pair query is answered positively (even if this does not reveal a $C_4$).
We shall refer to such an answer as a \emph{collision}. %\commentt{DR: better term?}

Let $\calG(K^{t-1})$ denote the subset of graphs  $G\in \Sup(\calD)$ such that %$\tau(G)$
$G$ is consistent with $K^{t-1}$. For a query $q^t$ (either a neighbor query or a pair query), let $\calG_c(K^{t-1},q^t)$ denote the subset of graphs  $G\in \calG(K^{t-1})$ for which the answer to $q^t$ according to %$\tau(G)$
$G$ results in a collision, and similarly define $\calG_{nc}(K^{t-1},q^t)$ for the case in which the answer is not a collision. % (henceforth, non-collision graphs).

\smallskip\noindent
\textsf{Subclaim}~
\emph{
Let $K^{t-1}$ be any knowledge graph where $t < n^{1/4}\alpha^{1/2}/c$, and let $q^t$ be a neighbor or pair query.
Then
$$ \frac{|\calG_{nc}(K^{t-1},q^t)|}{|\calG_{c}(K^{t-1},q^t)|} = \Omega(n^{1/4}\alpha^{1/2}) \;.$$
}

\smallskip
\noindent
\textsf{Proof:}~
In order to prove the subclaim we define an auxiliary bipartite graph $\calB(K^{t-1},q^t)$. Each node in $\calB(K^{t-1},q^t)$ corresponds to a graph $G \in \calG(K^{t-1})$.
% \commentt{DR: or should we say that its corresponding graph $G'$ is consistent with $K^{t-1}$, and maybe instead of $G'$ we need a notation like $\tau(G)$}.
Nodes on the left side correspond to graphs in $\calG_{c}(K^{t-1},q^t)$, and nodes on the right side to graphs in $\calG_{nc}(K^{t-1},q^t)$.
 There is an edge between two nodes corresponding to graphs $G_1$ and $G_2$, if $G_2$ can be obtained from $G_1$ by a certain ``swap'' modification (which will be explained precisely below). We shall show that the number of nodes on the left is smaller by a factor of $\Omega(n^{1/4}\alpha^{1/2})$ than the number of nodes on the left by appropriately lower bounding the degree of the former and upper bounding the degree of the latter.
 From this point on we use the shorthand notations $\calB$, $\calG$, $\calG_c$ and $\calG_{nc}$ for $\calB(K^{t-1},q^t)$, $\calG(K^{t-1})$, $\calG_c(K^{t-1},q^t)$ and $\calG_{nc}(K^{t-1},q^t)$, respectively.

% Consider first the case of a neighbor query $q^t = (x_{b,i},j)$ for $x_{b,i}\in X_b^{t-1}$, where $b\in \{1,2\}$.

%%%%%%%%%%%%%%%%%%%%%%%
\iffalse
For the sake of the exposition, suppose first for simplicity that graphs in $\calG$ have no parallel edges, so that $\tau(G)=G$ for every $G\in \calG$.  \commentt{DR: maybe say somewhere that if on parallel edges were created in $\calD$, then the entire argument would have been much simpler and explain where the difficulty is?}
\fi
%%%%%%%%%%%%%%%%%%%%%
For a graph $G\in \calG$ and an edge $e$ in $G$ we shall use the notation $\phi^{G}_X(e)$
and $\phi^G_Y(e)$ for its label with respect to its endpoint in $X$ and in $Y$, respectively.
Let $q^t$ be a neighbor query $nbr(v,i)$ where %(without loss of generality)
$v\in X^{t-1}$, 
% \commentt{DR: reminder to add the notation $nbr(\cdot)$ to the preliminaries}
and let $G_{c}$ be a graph in $\calG$ for which the answer to $q^t$ results in a collision. That is,  the answer $a^t$ is $(u,j)$ for some $u\in Y^{t-1}$.
The neighbors of node $G_{c}$ in $\calG_{nc}$ are nodes representing graphs that can be obtained from $G_{c}$ by swapping the edge $(v,u)$ with a different edge $(v, u')$ where $u'\notin Y^{t-1}$, and  %$G'_{c}$ is consistent with $K^{t-1}$.
such that the resulting graph belongs to $\calG$.
We prove that there are $\Omega(n^{1/2}\alpha)$ such graphs.

%    The swap modification will maintain the multiplicity of the edge $(v,u)$, and we start describing it for  the case that $(u,v)$ has multiplicity 1.
In order to modify the graph $G_c \in \calG_c$ to obtain a graph $G_{nc} \in \calG_{nc}$, we proceed as follows.
% Let $\overline{E}^{t-1}(v)$ denote the set of vertex in $Y^{t-1}$
    We choose a vertex %\commentt{DR: label?} 
    $u'\in Y\setminus Y^{t-1}$ (where the constraints on this choice are explained momentarily), and let $v'$ denote the $j$th neighbor of $u'$ in the graph $G_{c}$. Note that $v'$ is in $X$. Let $i'=\phi^{G_c}_X((v',u'))$. Now we swap the pair of  edges $(v,u)$
    and $(v',u')$   with the pair  $(v,u'), (v',u)$  while maintaining the original edge-labels. That is, $\phi^{G_{nc}}_X((v,u'))=i, \phi^{G_{nc}}_Y((v,u'))=j$ and $\phi^{G_{nc}}_X((v',u))=i', \phi^{G_{nc}}_Y((v,u'))=j'$.
    % (Note that this swap might result in multiple edges in the case that either $(v,u')$ or $(v',u)$ are already in  $G'_c$. )
    Since $u' \notin Y^{t-1}$, if $G_{nc}$ is in $\calG$, then it is in $\calG_{nc}$. In order to ensure that $G_{nc}\in\calG$, it is sufficient (and necessary), that there is no edge in $G_c$ between $v$ and $u'$ and no edge between $v'$ and $u$ and that
    $(v',u)\notin \overline{E}^{t-1}$.
    %$v'\notin X^{t-1}$
    %%%%%%%%%%%%%%%%%%%%%%%%%%
    \iffalse
    By the assumption that $u'$ is not in $Y^{t-1}$,
    and since the degrees of all involved vertices have not changed, and edges were only swapped,
    it holds that as long as $(v',u)\notin \overline{E}^{t-1}$, the resulting  graph is consistent with $K^{t-1}$.
    \fi
    %%%%%%%%%%%%%%%%

    Observe that the selection of $u'$ uniquely determines $v'$.
    There are $|Y^{t-1}| \leq t = O(n^{1/4}\alpha^{1/2})$ ``illegal'' choices of $u' \in Y^{t-1}$ and $2\alpha = O(n^{1/4}\alpha^{1/2})$ illegal choices of $u'$ that are neighbors of $v$. There are at most $t$ choices of $v'$ such that $(v',u)\in \overline{E}^{t-1}$, and $d$ choices of neighbors of $u$. For each of these $d+t$ choices of $v'$ there are at most $2\alpha$ neighbors $u'$ such that $v'$ may be the $j$th neighbor of $u'$. This rules out at most $2(d+t)\cdot\alpha$ choices of $u'$.
    Since $|Y|=|X|\alpha/d = n\alpha/(2d)$  \ForFuture{D: Be precise} and $d = n^{1/2}/c_2$, for a sufficiently large constant $c_2$, the number of choices of $u'$ for which $G_{nc} \in \calG$ is $\Omega(n^{1/2}\alpha)$.
    %%%%%%%%%%%%%%%%%%%%%%
    \iffalse
    Since there are $|Y|-|Y^{t-1}|=\Omega(n^{1/2}\alpha)$ options to choose $u'$ (which determines $v'$), and at most $t=O(n^{1/4}\alpha^{1/2})$ of them can result in an edge $(v',u)\in  \overline{E}^{t-1}$, there are $\Omega(n^{1/2}\alpha)$ graphs $G_{nc}$ that are neighbors of $G_{c}$ in $\calB$.
    \fi
    %%%%%%%%%%%%%%%%%%%%

    We now upper bound the number of neighbors of each node  in $\calB$ that corresponds to a graph $G_{nc} \in \calG_{nc}$.
    Assume that the answer to $q^t=nbr(v,i)$ in $G_{nc}$ is $a^t=(u',j')$. That is, $u'$ is the $i$th neighbor of $v$ in $G_{nc}$, and $v$ is the $j'$th neighbor of $u'$.
      By the definition of $\calB$, every neighbor of $G_{nc}$ in $\calB$ is a graph $G_{c}$ where the $i$th neighbor of $v$ is a vertex $u''\in Y^{t-1}$. Furthermore, by the definition of the swap process, if the edge $(v,u'')\in G_c$ was replaced by a swap process and resulted in an edge $(v,u')$ where $\phi^{G_{nc}}_Y((v,u'))=j'$, it must be the case that $\phi^{G_c}_{Y}((v,u''))=j'$ (since the swapping process  involves only two edges such that both have the same $\phi_Y$ value). Hence, for each choice of $u''\in Y^{t-1}$, the graph $G_c$ from which $G_{nc}$ was obtained by a swap process is uniquely defined:
      Let $v''$ denote the $j'$th neighbor of $u''$ in $G_c$.
      The graph $G_c$ is a graph where
      $(v,u'')$ is an edge with $\phi^{G_c}_{X}((v,u''))=i, \phi^{G_c}_Y((v,u''))=j'$ and $(v'', u')$ is an edge with $\phi^{G_c}_Y((v'', u'))=j'$ (and $\phi^{G_c}_X((v'', u'))=i''$ for some $i''$). Since there are at most $|Y^{t-1}|=O(n^{1/4}\alpha^{1/2})$ options to choose $u''$, it holds that $G_{nc}$ has $O(n^{1/4}\alpha^{1/2})$ neighbors in $\calB$ (it might be the case that some of the resulting graphs $G_c$ are not in
      % the support of $\calD$,
      $\calG$
      but the upper bound still holds). \ForFuture{D: doublecheck superscripts to $\phi$ (and other indices}

Since each node  corresponding to a graph in $\calG_c$ has $\Omega(n^{1/2}\alpha)$ neighbors and each node corresponding to a graph in $\calG_{nc}$ has $O(n^{1/4}\alpha^{1/2})$ neighbors, the subgraph follows for any neighbor query $nbr(v,i)$ for $v \in X^{t-1}$.

We next consider the case in which $q^t = nbr(u,j)$ for $u \in Y^{t-1}$. The swap operation and analysis are very similar to the case in which $q^t = nbr(v,i)$ for $v \in X^{t-1}$, and so we emphasize only the differences. Let  $G_c \in \calG_c$ be a graph for which $q^t$ is answered by $(v,i)$ where $v \in X^{t-1}$. To perform a swap operation and obtain a graph in $\calG_{nc}$ we select a vertex $v'\notin X^{t-1}$ such that $v'$ is not a neighbor of $u$ in $G_c$, and for the $i$th neighbor of $v'$, denoted $u'$ we have that
$u'$ is not a neighbor of $v$ in $G_c$, but
$(u',v)\notin \overline{E}^{t-1}$.  For each such legal choice of $v'$ we can perform a swap operation between the edges $(u,v)$ and $(u',v')$ while maintaining the edge labels as in the case that $q^t = nbr(v,i)$ for $v\in X^{t-1}$. The  constraints on $v'$ rule out at most $t+2d + (t+\alpha)\cdot 2d $ vertices, and so the number of legal choices of $v'$ is $\Omega(|X|) = \Omega(n)$. On the other hand, for any graph $G_{nc}\in \calG_{nc}$, for which the answer to $q^t = nbr(v,j)$ is $(v',i')$ for $v'\notin X^{t-1}$ there are at most $t$ choices of $v''\in X^{t-1}$ such that a swap operation between the edges $(v'',u)$ and $(v',u')$ resulted in the graph $G_{nc}$.
Hence, here the ratio between $|\calG_{nc}|$ and $|\calG_c|$ is $\Omega(n/t) = \Omega(n^{3/4}/\alpha^{1/2}) = \Omega(n^{1/4}\alpha^{1/2})$, as required.

 %     Before turning to the general case where there are parallel edges,
 Finally we turn to the case in which $q^t$ is a pair query between two vertices $v\in X^{t-1}$ and $u\in Y^{t-1}$. Consider a graph $G_c \in \calG_c$ for which the answer is positive.
 % and let $i$ be the label of the edge from its endpoint  $v$, and $j$ from its endpoint $u$.
 We can perform a swap with any edge $(v',u')$ in $G_c$ such that
 $(v',u')\notin E(K^{t-1})$ and  in addition, $(v',u)\notin E(G_c)$, $(v,u')\notin E(G_c)$ and $(v',u), (v,u')\notin \overline{E}^{t-1}$. The number of edges in $G_c$ on which a swap cannot performed with $(u,v)$ is hence upper bounded by $t+ (d+t)\cdot \alpha + d\cdot (t+\alpha) = O(d\alpha)$. This leaves $\Omega(n\alpha)$ edges with which a swap can be performed (each resulting in a different graph in $\calG_{nc}$). On the other hand, for a graph $G_{nc}\in \calG_{nc}$, for which there is no edge between $u$ and $v$, the number of graphs in $\calG_c$ for which a swap operation resulted in $G_{nc}$ is at most $d\alpha$ (the number of pairs of vertices $(u',v')$ such that $(u,v'),(u',v)\in E(G_{nc})$). Hence the ratio between $|\calG_{nc}|$ and $|\calG_c|$ is $\Omega((n\alpha)/(d\alpha) = \Omega(n/d) = \Omega(n^{1/2})= \Omega(n^{1/4}\alpha^{1/2})$.
 %%%
% Observe that the above argument for the case that $q^t$ is a neighbor query is applicable essentially as is also to the case that $q^t$ is a pair query on a pair of vertices $(u,v)$ (where graphs in $\calG_c$ have an edge between the two vertices and graphs in $\calG_{nc}$ do not).
%%%%%%%%%%%%%%%%%%%%%%%%%%%%%%%%%%%%%%%%
\iffalse
It remains to address the general  case where graphs in $\calG$ have parallel edges. Once again, consider the case that $q^t$ is a neighbor query on a vertex $v\in X^{t-1}$ (here too the case in which it is a pair query is dealt with essentially the same). Consider any graph $G_c \in \calG_c$ such that  the neighbor query $q^t = (v,i)$  is answered in $\tau(G_c)$ by $(u,j)$ for some $u \in Y^{t-1}$. Denoting by $\ell$ the multiplicity of the edge $(u,v)$ in $G_c$, let $i_1,\dots,i_\ell$ and $j_1,\dots,j_\ell$ be such that for each $r \in [\ell]$, there is an edge $e_r$ in $G_c$ between $u$ and $v$ such that $\phi^{G_c}_X(e_r)=i_r$ and $\phi^{G_c}_Y(e_r)=j_r$.
      In this case all these edges are swapped simultaneously as described next (by generalizing the swap operation for the case when there are no parallel edges).
 %     Let $i=i_1, \ldots, i_l$ be the indices of edges from $v$ to $u$ from the perspective of $v$, and $j=j_1, \ldots, j_{\ell}$ the indices of these edges from the perspective of $u$. To swap the hyper-edge To swap the hyper-edge $(v,u)$,

 We choose a vertex $u'\in Y\setminus Y^{t-1}$  and consider its $j_1,\dots,j_\ell$ incident edges in $G_c$, denoted $e'_1,\dots,e'_\ell$, and the corresponding endpoints,
$v_1, \ldots, v_{\ell}$ (where these vertices are not necessarily the same nor necessarily distinct).
 In order to obtain $G_{nc}$ for this choice of $u'$, we first remove the $\ell$ parallel edges between $v$ and $u$ as well as the edges %between $u$ and $v_{j_1},\dots,v_{j_\ell}$.
 $e'_1,\dots,e'_\ell$ that are incident to $u'$.
 We then  add $\ell$ parallel edges $e_1,\dots,e_\ell$ between $v$ and $u'$
 % an hyper-edge with the  multiplicity  $\ell$ between $v$ to $u'$,
 such that $\phi^{G_{nc}}_X(e_r)= i_r$ and $\phi^{G_{nc}}_Y(e_r)= j_r$ for each $r \in [\ell]$. Finally, we  add edges $e''_1,\dots,e''_\ell$ where $e''_r$ for $r\in [\ell]$ is between $v_{r}$ and $u$, $\phi^{G_{nc}}_X(e''_r) = \phi^{G_c}_X(e'_r)$, and $\phi^{G_{nc}}_Y(e''_r) = j_r$.
 % $(v_{j_1}, u), \ldots, (v_{j_{\ell}}, u)$, such that for every $r\in [1,\ell]$, it holds that  $\phi_X((v_{j_k},u))=\phi_X((u', v_{j_{k}}))$ and $\phi_Y((v_{j_k},u))=j_k$.


  Here too each choice of $u'\in Y\setminus Y^{t-1}$ uniquely determines the vertices $v_1, \ldots v_{\ell}$, and unless $(v_{r},u)\in \overline{E}^{t-1}$, the resulting graph $G_{nc}$ after the swap is consistent with $K^{t-1}$. Hence,  $G_c$ has $\Omega(n^{1/2}\alpha)$ neighbors in $\calB$.
   \commentt{DR: I have to finish from here on, but am a bit unsure about the last statement because of the constraint on non-edges and the possibility of high-multiplicity}
    On the other hand,  consider any graph $G_{nc} \in \calG_{nc}$.
    For such a graph, the query $q^t = (v,i)$ is answered according to $\tau(G_{nc})$ by $(u',j')$ for $u' \notin Y^{t-1}$.
    Let $e'_1,\dots,e'_{\ell'}$ be the parallel edges between $v$
    and $u'$ in $G_{nc}$ where $\phi^{G_{nc}}_X(e'_r)=i'_r$ and
    $\phi^{G_{nc}}_Y(e'_r)=j'_r$.
    A swap operation on a graph $G_{c} \in \calG_c$ results in $G_{nc}$, only if there exists $u''\in Y^{t-1}$ such that
    there are $\ell'$ parallel edges $e''_1,\dots,e''_{\ell'}$ between $v$ and $u''$ in $G_c$ where $\phi^{G_C}_Y(e''_r)=j_r$
    %and $\phi^{G_c}_X(e''_r)=i_r$.
    Therefore, for every  $u''\in Y^{t-1}$, there  is at most  one graph $G_c$ for which the swap operation results in $G_{nc}$, implying that for each $G_{nc}\in \calG_{nc}$, its degree in $\calB$ is $O(n^{1/4}\alpha^{1/2})$. This concludes the proof of the subclaim.
\fi
%%%%%%%%%%%%%%%%%%%
  {\hspace*{\fill}$\Diamond$\par}
% \end{proof}

\smallskip
By the above, the probability that the algorithm detects a $C_4$ following any one of its at most
$n^{1/4}\alpha^{1/2}/c$ queries is upper bounded by a small constant (e.g., $1/10$).
\end{proof}


%%%%%%%%%%%%%%%%%%%%%%%%%%%%%%%%%
\iffalse
\begin{claim}
    Let $q_t$ be a neighbor query $nbr(v,i)$ for some $v\in X$.
    The number of graphs in the support of $\calD$ that correspond to a collision is smaller by a factor of $\Omega(n^{1/4}\alpha^{1/2})$ than the number of graphs that do not correspond to a collision (henceforth, non-collision graphs).
\end{claim}
\begin{proof}
    We define an  auxiliary bipartite  graphs $\calB$, where each node in $\calB$ corresponds to a graph in the support of $\calD$ that is consistent with $K^{t-1}$.
   Graphs in $\calD$ for which the answer $a^t$ to $q^t$ corresponds to a collision, belong to the left side of $\calB$, denoted $\calB_{c}$, and the others to the right side, denoted $\calB_{nc}$.


    There is an edge between two nodes in (different sides of) $\calB$, if the graphs they correspond to can be obtained from one another by a ``swap'' modification, which we define below.
    We shall prove that nodes  on the left side have degree $\Omega(n^{1/2}\alpha)$, while nodes on the right side have degree $O(n^{1/4}\alpha^{1/2})$. Since $\calB$ is bipartite, this implies that $|\calB_{nc}|$ is greater by a factor of $\Omega(n^{1/4}\alpha^{1/2})$ than $\calB_{c}$. By the definition of $\calB$ this will conclude the proof.

    Assume first that the query $q^t$ is a neighbor query $(v,i)$, and let $G_{c} \in \calD$ be a graph whose answer to the query results in a collision. That is, assume that the answer $a^t$ is $(u,j)$ for some $u\in K^t$. The neighbors of node $G_{c}$ in $\calB_{nc}$ are nodes that represent graphs that can be obtained from $G_{c}$ by swapping the (hyper-)edge $(v,u)$ in a different (hyper-)edge $(v, u')$ where $u'\notin K^{t-1}$, and  $G'_{c}$ is consistent with $K^{t-1}$. We prove that there are $\Omega(n^{1/2}\alpha)$ such graphs.

    The swap modification will maintain the multiplicity of the edge $(v,u)$, and we start describing it for  the case that $(u,v)$ has multiplicity 1.  In order to modify the graph $G_c$ to get a new graph $G_{nc}$ that is consistent with $K_{t-1}$, we proceed as follows.
    We choose a label  $u'\in Y\setminus V(K_t)$ u.a.r., and let $v'$ denote the $j$th neighbor of $u'$ in the graph $G_{c}$. Note that $v'$ is in $X$. Let $i'$ denote $\phi_X((v',u'))$. Now we swap the pair of  edges $(v,u)$
    and $(v',u')$   with the pair  $(v,u'), (v',u)$  while maintaining the original labels. That is, $\phi_X((v,u'))=i, \phi_Y((v,u'))=j$ and $\phi_X((v',u))=i', \phi_Y((v,u'))=j$.
    % (Note that this swap might result in multiple edges in the case that either $(v,u')$ or $(v',u)$ are already in  $G'_c$. )
    By the assumption that $u'$ is not in $K^{t-1},$
    and since the degree of all involved vertices had not changed, and edges were only swapped,
    it holds that as long as $(v',u)\notin \bar{E}^{t-1}$, the resulting  graph is consistent with $K^{t-1}$.
    Also, observe that the selection of $u'$ uniquely defines $v'$.
    Since there are $|Y|-|V(K^{t-1})|=\Omega(n^{1/2}\alpha)$ options to choose $v'$ and at most $t=O(n^{1/4}\alpha^{1/2})$ of them can result in an edge $(v',u)\in  \bar{E}^{t-1}$, there are $\Omega(n^{1/2}\alpha)$ graphs $G_{nc}$ that are neighbors of $G_{c}$ in $\calB$.

    We now upper bound the number of neighbors of each $G_{nc}$ in $\calB$.
    Assume that the answer to $a_t=nbr(v,i)$ in $G_{nc}$ is $(u',j')$. That is, $u'$ is the $i$th neighbor of $v$ in $G_{nc}$, and $v$ is the $(j')$th neighbor of $u'$.
      By the definition of $\calB$, every neighbor of $G_{nc}$ in $\calB$ is a graph $G_{c}$ where the $i$th neighbor of $v$ is a vertex $u''\in K^{t-1}$. Furthermore, by the definition of the swap process, ifthe edge $(v,u'')\in G_c$ was replaced by a swap process and resulted in edge $(v,u')$ where $\phi_Y((v,u'))=j'$, it must be the case that $\phi_{Y}((v,u''))=j'$ (since the swapping process  involves only two edges such that both have the same $\phi_Y$ value). Hence, for each choice of $u''\in K^{t-1}$, the graph $G_c$ from which $G_{nc}$ was obtained by a swap process is uniquely defined:
      Let $v''$ denote the $(j')$ neighbor of $u''$ in $G_c$.
      The graph $G_c$ is a graph where
      $(v,u'')$ is an edge with $\phi_{X}((v,u''))=i, \phi_Y((v,u''))=j'$ and $(v'', u')$ is an edge with $\phi_Y((v'', u'))=j'$ (and $\phi_X((v'', u'))=i''$ for some $i''$). Since there are at most $|V(K^{t-1})|=O(n^{1/4}\alpha^{1/2})$ options to choose $u''$, it holds that $G_{nc}$ has at most $=O(n^{1/4}\alpha^{1/2})$ neighbors in $\calB$ (it might be the case that some of the resulting graphs $G_c$ are not in the support of $\calD$, but the upper bound still holds).

      It remains to address the case where the answer to $a_t$ is an edge $(v,u)$ with multiplicity $\ell>1$.
      % In this case all the copies of the edge $(v,u)$ are swapped simultaneously as follows. Let $i=i_1, \ldots, i_l$ be the indices of edges from $v$ to $u$ from the perspective of $v$, and $j=j_1, \ldots, j_{\ell}$ the indices of these edges from the perspective of $u$. To swap the hyper-edge  $(v,u)$, we choose a vertex $u'\in Y\setminus K^{t-1}$ u.a.r., and look at its $j_1$ to $j_l$th neighbors,  denoted  $v_{j_1}, \ldots, v_{j_{\ell}}$, respectively. We now remove the hyper-edge $(v,u)$ as well as the edges $(u, v_{j_1}), \ldots, (u,v_{j_{\ell}})$.
      % % We then  add an hyper-edge with the  multiplicity  $\ell$ between $v$ to $u'$, such that $\phi_Y((v,u'))=\{j_1, \ldots, j_{\ell}\}$, and add edges $(v_{j_1}, u)$
      In this case all the copies of the edge $(v,u)$ are swapped simultaneously as follows. Let $i=i_1, \ldots, i_l$ be the indices of edges from $v$ to $u$ from the perspective of $v$, and $j=j_1, \ldots, j_{\ell}$ the indices of these edges from the perspective of $u$. To swap the hyper-edge To swap the hyper-edge $(v,u)$, we choose a vertex $u'\in Y\setminus V(K^{t-1})$ u.a.r., and look at its $j_1$ to $j_l$th neighbors,  denoted  $v_{j_1}, \ldots, v_{j_{\ell}}$. We now remove the hyper-edge $(v,u)$ as well as the edges $(u, v_{j_1}), \ldots, (u,v_{j_{\ell}})$. We then  add an hyper-edge with the  multiplicity  $\ell$ between $v$ to $u'$, such that $\phi_Y((v,u'))=\{j_1, \ldots, j_{\ell}\}$, and add edges $(v_{j_1}, u), \ldots, (v_{j_{\ell}}, u)$, such that for every $k\in [1,\ell]$, it holds that  $\phi_X((v_{j_k},u))=\phi_X((u', v_{j_{k}}))$ and $\phi_Y((v_{j_k},u))=j_k$.
      Each choice of $u'\in Y\setminus V(K^{t-1})$ uniquely defines the set $v_{j_1}, \ldots v_{j_\ell}$, and unless $(v_{j_k},u)\in \bar{E}^{t-1}$, the resulting graph $G_{nc}$ after the swap is consistent with $K^{t-1}$. Hence,  $G_c$ has $\Omega(n^{1/2}\alpha^{1/4})$ neighbors in $\calB$.
    On the other hand, similarly to what we had before, a graph $G_{c}^*$ can be swapped to $G_{nc}$ above, only if there exists $u''\in V(K^{t-1})$ such that    $(v,u'')$ in $E(G_{c})$ has multiplicity $\ell$ and $\phi_Y((v,u''))=\{j_1, \ldots, j_{\ell}\}$. Therefore, for every  $u''\in V(K^{t-1})$, there  is at most  one graph $G_c^*$ that can be swapped to $G_{nc}$, implying that for each $G_c\in \calB$, its degree in $\calB$ is $O(n^{1/4}\alpha^{1/2})$. This concludes the proof of the subclaim.

%%%%%%%%%%%%

%      \commentt{T: I proved the claim on graphs in $\calD$ rather than in $\calD'$ - I'm not sure if that's sufficient}


\end{proof}
\fi



\bigskip
Theorem~\ref{thm:C4-gen-alpha-lb} follows by combining Lemmas~\ref{lem:D-C4-far} and~\ref{lem:D-C4-lb}.
\ForFuture{D: or spell it out}

%%%%%%%%%%%%%%%%%%%%%%%%%%%%%%%%%%%%
\iffalse
\newpage
\section{From writeup}

\commentt{DR: Base on the following sketch (from the summaries file)}


First we recalled that we have an $\Omega(n^{1/4})$ two-sided error lower bound for constant $\alpha$. I think it is written somewhere, but let's recall quickly. In the FREE instances we have a bipartite graph with $\Theta(n)$ vertices on one side (the left) and $\Theta(n^{1/2})$ on the other. Every vertex on the left has two unique neighbors on the right. In the FAR instances we pair the vertices on the left, and each pair has the same two neighbors on the right. I guess on the right we need to have less vertices, so that we'll also have a small independent set of ``leftover'' vertices. Can consider other (more elegant?) variants, which we may have though of in the past.

Then we recalled the construction for the case that $\alpha = \Theta(n^{1/2})$, for a one-sided error lower bound of $\Omega(n^{1/2})$ (which is $\Omega(n^{1/4} \alpha^{1/2})$). Here (only have FAR instances), on each side we have $\Theta(n)$ vertices with $\Theta(n^{1/2})$ perfect matchings between them. As we noted before, this creates parallel edges, so we need to think whether/how to deal with this. The easy part here is to show the lower bound on the number of queries necessary to see a C4  with constant probability, and the more subtle part (which I think maybe we didn't do in full detail in the past) is showing that whp we indeed get a graph that is far from C4-free. I think in the end (of this meeting) we didn't quite get this to work for $\alpha = \Theta(n^{1/2})$ and constant distance, but only for slightly smaller $\alpha$, and since  both parts of the argument are quite similar for general $\alpha$ (or at least, $\alpha \leq n^{1/2}/\log n$ or so), will turn to that.

The construction for general $\alpha$ still has on one side $\Theta(n)$ vertices (let's say $n$ for simplicity), but on the other side it has $n^{1/2}\alpha$ vertices. The vertices on the lhs have degree $\alpha$, and those on the right have degree $n^{1/2}$.
The edges are selected by picking a random matching between entries of an $n\times \alpha$ table (corresponding to the ``neighbor/edge ports'' of the lhs vertices) and the entries of an $n^{1/2}\alpha \times n^{1/2}$ table (corresponding to the ``neighbor/edge ports of the rhs vertices). Here parallel edges will appear even more, but we continue glossing over this.

To show that $\Omega(n^{1/4} \alpha^{1/2})$ queries are necessary to obtain a C4, we actually show that these many queries are necessary for viewing any cycle. Consider any knowledge graph that is obtained after $n^{1/4} \alpha^{1/2}/c$ queries (for some sufficiently large constant $c$). This graph contains at most $2 n^{1/4} \alpha^{1/2}/c$ vertices. So now, if we condition on this knowledge graph and bound the probability that a neighbor query $(v,i)$ will be answered by some vertex $u$ that is in the knowledge graph, we get that it is at most $2 n^{1/4} \alpha^{1/2}/c$ divided by $n^{1/2}\alpha$, which is $2/(cn^{1/4}\alpha^{1/2})$ (we need to be a bit more precise here -- the conditioning on edges already in the graph only reduces the probability of hitting a vertex in the knowledge graph, but there are also non-edges, and these maybe have an opposite effect). For pair queries it looks like it will be something like $\alpha /(n^{1/2}\alpha) = 1/n^{1/2}$.

Let's turn to showing that we get whp a graph that has $\Omega(m)$ disjoint C4s. The first part of the argument is showing that for every edge, the expected number of C4s that it participates in is $\alpha$.
Consider an edge $(u,v)$. For any choice of sets of neighbors $\Gamma(u)$ and $\Gamma(v)$, and for each pair $(u',v') \in \Gamma(v) \times \Gamma(u)$, the probability that we have an edge between them is (roughly) $\alpha \times 1/(n^{1/2}\alpha) = 1/n^{1/2}$. Since we have $n^{1/2} \cdot \alpha$ such pairs, we get  $\alpha$ in expectation. Hence, for any factor of $2^i$, the fraction of edges that participate in more than $2^i\cdot \alpha$ C4s, is at most $1/2^i$, where the maximum number of C4's that an edge can participate in, is $n^{1/2}\cdot \alpha$. We did a simple bucketing argument, with $O(\log n)$ buckets, by which we got that after removing all edges that reside on more than $c' \alpha$ C4s, we removed at most $\log n \cdot n\alpha \cdot \alpha$ C4s, which for $\alpha$ sufficiently smaller than $(n/\log n)^{1/2}$, is much smaller than $n^2$.

Now we use the fact that the expected number of C4s is $n^4 \cdot (1/n^{1/2})^4 = n^2$. I guess here we need to show that whp we won't get significantly less than this expected value (expectation is not enough, and the C4s are not independent of each other, so maybe need Chebishev).

So we should be left with $\Omega(n^2)$ C4s, where each remaining edge resides on at most $c'\alpha$ (maybe $c'$ is simply $ 2$), so if we run a greedy process for $n\alpha/c''$ iterations, where in each iteration we  assign a ``free'' C4 to an edge (and mark all C4s it participates in as ``un-free''), we'll ``use up'' at most $(n \alpha/c'') \cdot c'\alpha$ C4s, which is less than $(c'/c) n^2$.

Other than doing things more carefully and precisely, there is a question as to whether the $\log n$ factor is really necessary.
\fi
%%%%%%%%%%%%%%%%%%%%%%%%%%%


